%%=============================================================================
%% De inhoud van je projectvoorstel
%%=============================================================================

%---------- 1. Context en probleemstelling ------------------------------------
\section{Context en probleemstelling}%
\label{vst:context}
Tijdens mijn stage bij Vesalius.ai werk ik mee aan een digitaal platform dat
artsen in ziekenhuizen ondersteunt bij consultaties en de administratieve
verwerking ervan. Momenteel kunnen artsen en organisaties de inhoud en opmaak
van de meeste automatisch verstuurde e-mails, zoals afspraakbevestigingen,
herinneringen en annulaties, niet zelf aanpassen, waardoor de communicatie
onvoldoende flexibel is om aan hun specifieke behoeften te voldoen.

%---------- 2. Doelstelling ---------------------------------------------------
\section{Doelstelling}%
\label{vst:doelstelling}
Het doel is een werkende proof of concept (PoC) te ontwikkelen van een visuele
e-maileditor waarmee gebruikers de tekst en opmaak van drie vooraf bepaalde
e-mailtypes kunnen aanpassen, opslaan, vooraf bekijken en als testmail
versturen. Het project is geslaagd wanneer de aangepaste e-mails correct worden
gerenderd en de weergave ervan in gangbare e-mailclients, waaronder Gmail,
Outlook.com en Apple Mail, systematisch werd getest en gedocumenteerd.

%---------- 3. Technologieën --------------------------------------------------
\section{Technologieën}%
\label{vst:technologieen}
De frontend wordt ontwikkeld met Angular 21 en TypeScript, met Quill 2 als
basis voor de visuele editor; voor de mock-API worden Node.js, TypeScript en
Express gebruikt. Voor de verwerking en beveiliging van e-mails worden onder
meer Handlebars, \texttt{sanitize-html}, \texttt{juice} en
\texttt{html-to-text} ingezet; mockgegevens, geautomatiseerde tests met Vitest
en Playwright en Docker met Mailpit ondersteunen de ontwikkeling en validatie.
De oplossing wordt ontwikkeld in Bitbucket, conform de voorkeur van het
bedrijf, en maakt uitsluitend gebruik van fictieve gegevens, zodat geen
reële patiëntgegevens worden verwerkt; AI wordt niet ingezet, aangezien het
project zich richt op deterministische tekstbewerking en e-mailrendering.
Bij de ontwikkeling wordt aandacht besteed aan veilige HTML-verwerking,
placeholders en de relevante principes van de AVG.

%---------- 4. Afbakening -----------------------------------------------------
\section{Afbakening}%
\label{vst:afbakening}
Mijn dagelijkse stagewerk bestaat uit het analyseren en oplossen van bugs en
het uitvoeren van andere ontwikkeltaken binnen het bestaande Vesalius-platform.
Het graduaatsproject is hiervan afzonderlijk afgebakend en omvat de analyse,
het ontwerp, de realisatie en het testen van een zelfstandige PoC met één
visuele editor, opslag van e-mailtemplates, ondersteuning voor placeholders,
een voorbeeldweergave en het versturen van testmails voor drie e-mailtypes.
De oplossing werkt met fictieve gegevens en een tijdelijke testopstelling;
integratie in het bestaande platform, productiegebruik, Keycloak-e-mails,
productieplanning en een volledig configureerbare layout-builder vallen buiten
de scope.

%---------- 5. Verwachte meerwaarde -------------------------------------------
\section{Verwachte meerwaarde}%
\label{vst:meerwaarde}
De PoC toont aan in welke mate artsen en organisaties meer controle kunnen
krijgen over de inhoud en opmaak van hun e-mailcommunicatie zonder dat de
bestaande e-mailstructuur volledig moet worden vervangen. Het bedrijf ontvangt
een werkende demonstratie, een testopstelling met herbruikbare mockgegevens,
geautomatiseerde tests en een overzicht van de resultaten per e-mailclient.
Daarnaast wordt een onderbouwd integratievoorstel opgeleverd met aandacht voor
architectuur, beveiliging en mogelijke hergebruiksmogelijkheden binnen het
bestaande Vesalius-platform.

%---------- 6. Planning en werkbelasting --------------------------------------

\section{Planning en werkbelasting}%
\label{vst:planning}
Het graduaatsproject omvat ongeveer 100 uur, gespreid over tien wekelijkse
sprints van telkens één gefocuste werkdag, aangevuld met maximaal 20 uur
buffertijd voor eventuele correcties en de voorbereiding van de eindpresentatie.
De planning en voortgang worden regelmatig met de bedrijfsmentor afgestemd;
per sprint wordt een werkend resultaat nagestreefd en gedemonstreerd.

\begin{table}[htbp]
  \centering
  \begin{tabular}{p{0.25\textwidth}p{0.57\textwidth}r}
    \toprule
    \textbf{Fase} & \textbf{Wat lever je op} & \textbf{Uren} \\
    \midrule
    Analyse en vereisten
      & Scope, functionele vereisten en testcriteria
      & 12 \\
    Technologiekeuze en opzet
      & Gemotiveerde technologiekeuzes en werkende ontwikkelomgeving
      & 10 \\
    Ontwerp
      & Architectuur, editorconcept, templateopbouw en testmatrix
      & 13 \\
    Realisatie
      & Werkende editor, templateopslag, renderpipeline, preview en testmails
      & 40 \\
    Testen en validatie
      & Geautomatiseerde tests en resultaten van e-mailclienttests
      & 15 \\
    Rapport, poster en demo
      & Eindrapport, integratievoorstel en demonstratie
      & 10 \\
    \midrule
    \textbf{Totaal} & & \textbf{100} \\
    \bottomrule
  \end{tabular}
  \caption[Planning en werkbelasting]{\label{tab:vst-planning}%
    Indicatieve verdeling van de 100 projecturen.}
\end{table}

%---------- 7. Risico's -------------------------------------------------------

\section{Risico's}%
\label{vst:risicos}
Omdat het project grotendeels zelfstandig en met een beperkte wekelijkse
werkbelasting wordt uitgevoerd, worden de taken per sprint afgebakend en de
voortgang regelmatig met de bedrijfsmentor besproken. Verschillen tussen de
mock-opstelling en de bestaande Laravel-omgeving worden gedocumenteerd in een
integratievoorstel, zodat de PoC geen volledige integratie vereist. Indien de
gekozen editor onvoldoende mogelijkheden biedt of e-mailclients de uitvoer
anders weergeven dan verwacht, worden alternatieve editoropties tijdig
geëvalueerd en worden de belangrijkste compatibiliteitsproblemen prioritair
opgelost binnen de voorziene scope en buffertijd.
